\subsection{Đề 2}
\Opensolutionfile{ans}[ans/ans-2-B11-De2]
\caulc
%%==========Câu 1
\begin{ex}%[2D4N1-1]%[2015_K12 Huyên Nguyễn]
	Cho $F(x)$ là nguyên hàm của hàm số $f(x)=\sin \left(\dfrac{\pi}{2}-x\right)$. Khi đó $F'(x)$ bằng
	\choice
	{$\cos \left(\dfrac{\pi}{2}-x\right)$}
	{$\sin x$}
	{$-\cos \left(\dfrac{\pi}{2}-x\right)$}
	{\True $\cos x$}
	\loigiai{
		Vì $F(x)$ là một nguyên hàm của hàm số $f(x)$ nên $F'(x)=f(x)=\sin \left(\dfrac{\pi}{2}-x\right)=\cos x$.
	}
\end{ex}

%%==========Câu 2
\begin{ex}%[2D4N1-2]%[2015_K12 Huyên Nguyễn]
	Họ nguyên hàm của hàm số $f(x)=2x+\dfrac{3}{x^2}$ là
	\choice
	{\True $x^2-\dfrac{3}{x}+C$}
	{$x^2+\dfrac{3}{x}+C$}
	{$x^2+3\ln{x^2}+C$}
	{$x^2+\dfrac{3}{2}\ln|x^2|+C$}
	\loigiai{
		$\displaystyle\int\left(2x+\dfrac{3}{x^2}\right)\mathrm{\,d}x=x^2-\dfrac{3}{x}+C$.
	}
\end{ex}

%%==========Câu 3
\begin{ex}%[2D4N1-2]%[2015_K12 Huyên Nguyễn]
	Cho $F(x)$ là một nguyên hàm của hàm số $f(x)=\dfrac{1}{x}$ và thỏa mãn $F\left(\mathrm{e}^2\right)=3$. Khi đó $F(x)$ bằng
	\choice
	{$\ln |x|+5$}
	{\True $\ln |x|+1$}
	{$\ln |x|+3$}
	{$\ln |x|-1$}
	\loigiai{
		Ta có $F(x)=\displaystyle\int\dfrac{1}{x}\mathrm{\,d}x=\ln|x|+C$.\\
		Mà $F\left(\mathrm{e}^2\right)=3\Leftrightarrow \ln|\mathrm{e}^2|+C=3\Leftrightarrow C=1$.\\
		Vậy $F(x)=\ln |x|+1$.
	}
\end{ex}

%%==========Câu 4
\begin{ex}%[2D4N1-3]%[2015_K12 Huyên Nguyễn]
	$\displaystyle\int \dfrac{1}{1+\cos 2x} \mathrm{\,d}x$ bằng
	\choice
	{\True $\dfrac{1}{2}\tan x+C$}
	{$2\tan x+C$}
	{$\dfrac{1}{2}\tan 2x+C$}
	{$2\tan 2x+C$}
	\loigiai{
		Ta có $\displaystyle\int \dfrac{1}{1+\cos 2x} \mathrm{\,d}x=\displaystyle\int \dfrac{1}{2\cos^2x} \mathrm{\,d}x=\dfrac{1}{2}\tan x+C$.
	}
\end{ex}

%%==========Câu 5
\begin{ex}%[2D4H1-4]%[2015_K12 Huyên Nguyễn]
	$\displaystyle\int\left(\dfrac{1}{\mathrm{e}^{3x}}+1\right)\mathrm{\,d}x$ bằng
	\choice
	{\True $-\dfrac{1}{3 \mathrm{e}^{3x}}+x+C$}
	{$\ln \left(\mathrm{e}^{3x}\right)+x+C$}
	{$-\dfrac{3}{\mathrm{e}^{3x}}+x+C$}
	{$\dfrac{1}{3}\ln \left(\mathrm{e}^{3x}\right)+x+C$}
	\loigiai{
		Ta có $\displaystyle\int\left(\dfrac{1}{\mathrm{e}^{3x}}+1\right)\mathrm{\,d}x=\int \left( e^{-3x}+1\right)\mathrm{\,d}x=-\dfrac{1}{3} \mathrm{e}^{-3x}+x+C$.
	}
\end{ex}

%%==========Câu 6
\begin{ex}%[2D4H1-4]%[2015_K12 Huyên Nguyễn]
	Cho $F(x)$ là một nguyên hàm của hàm số $f(x)=\mathrm{e}^{2x}$ và thỏa mãn $F(\ln 2)=5$. Khi đó $F(x)$ bằng
	\choice
	{$\dfrac{1}{2}\mathrm{e}^{2x}+4$}
	{$2 \mathrm{e}^{2x}+1$}
	{\True $\dfrac{1}{2}\mathrm{e}^{2x}+3$}
	{$2 \mathrm{e}^{2x}-3$}
	\loigiai{
		Ta có $F(x)=\displaystyle\int\mathrm{e}^{2x}\mathrm{\,d}x=\dfrac{1}{2}\mathrm{e}^{2x}+C$.\\
		Do $F(\ln 2)=5\Leftrightarrow 2+C=5\Leftrightarrow C=3$.\\
		Vậy $F(x)=\dfrac{1}{2}\mathrm{e}^{2x}+3$.
	}
\end{ex}

%%=====Câu 7
\begin{ex}%[2D4H1-2]%[2015_K12 Huyên Nguyễn]
	$\displaystyle\int\dfrac{1-x}{1+\sqrt{x}}\mathrm{\,d}x$ bằng
	\choice
	{$x-\dfrac{1}{\sqrt{x}}+C$}
	{$x-\dfrac{1}{2\sqrt{x}}+C$}
	{$x-\dfrac{3}{2}x\sqrt{x}+C$}
	{\True $x-\dfrac{2}{3}x\sqrt{x}+C$}
	\loigiai{
		$\displaystyle\int\dfrac{1-x}{1+\sqrt{x}}\mathrm{\,d}x=\displaystyle\int\dfrac{\left(1-\sqrt{x}\right)\left(1+\sqrt{x}\right)}{1+\sqrt{x}}\mathrm{\,d}x=\displaystyle\int\left(1-\sqrt{x}\right)\mathrm{\,d}x=x-\dfrac{2}{3}x\sqrt{x}+C$.
	}
\end{ex}

%%=====Câu 8
\begin{ex}%[2D4H1-2]%[2015_K12 Huyên Nguyễn]
	Một nguyên hàm của $f(x)=\dfrac{1}{\sqrt{x+9}-\sqrt{x}}$ là 
	\choice
	{$\dfrac{2}{3\left[\sqrt{(x+9)^3}-\sqrt{x^3}\right]}$}
	{$\dfrac{2}{27}\left[\sqrt{(x+9)^3}-\sqrt{x^3}\right]$}
	{\True $\dfrac{2}{27}\left[\sqrt{(x+9)^3}+\sqrt{x^3}\right]$}
	{$\dfrac{2}{3\left[\sqrt{(x+9)^3}+\sqrt{x^3}\right]}$}
	\loigiai{Vì $\left(\sqrt{x+9}-\sqrt{x}\right)\left(\sqrt{x+9}+\sqrt{x}\right)=9$ nên $f(x)=\dfrac{1}{\sqrt{x+9}-\sqrt{x}}=\dfrac{\sqrt{x+9}+\sqrt{x}}{9}$.\\
		Vậy một nguyên hàm của $f(x)$ là $\dfrac{2}{27}\left[\sqrt{(x+9)^3}+\sqrt{x^3}\right]$.
	}
\end{ex}

%%=====Câu 9
\begin{ex}%[2D4H1-3]%[2015_K12 Huyên Nguyễn]
	Họ nguyên hàm của hàm số $f(x)=\dfrac{3}{\cos^2 2x}$ là
	\choice
	{\True $\dfrac{3}{2}(1+\tan 2x)+C$}
	{$6(1+\tan 2x)+C$}
	{$6\left(1+\tan^2 2x\right)+C$}
	{$\dfrac{3}{2}\left(1+\tan^2 2x\right)+C$}
	\loigiai{
		Ta có $\displaystyle\int \dfrac{3}{\cos^2 2x}\mathrm{\,d}x=\dfrac{3}{2}(1+\tan 2x)+C$.
	}
\end{ex}

%%=====Câu 10
\begin{ex}%[2D4H1-3]%[2015_K12 Huyên Nguyễn]
	Cho $F(x)$ 1à một nguyên hàm $f(x)=\sin x-\cos x$ và thỏa $F\left(\dfrac{\pi}{4}\right)=0$. Khi đó $F(x)$ bằng
	\choice
	{\True $-\cos x-\sin x+\sqrt{2}$}
	{$-\cos x-\sin x+\dfrac{\sqrt{2}}{2}$}
	{$\cos x-\sin x+\sqrt{2}$}
	{$\cos x-\sin x+\dfrac{\sqrt{2}}{2}$}
	\loigiai{
		Ta có $F(x)=\displaystyle\int\left(\sin x-\cos x\right)\mathrm{\,d}x=-\cos x-\sin x+C$.\\
		Mà $F\left(\dfrac{\pi}{4}\right)=0\Leftrightarrow -\sqrt{2}+C=0\Leftrightarrow C=\sqrt{2}$.\\
		Vậy $F(x)=-\cos x-\sin x+\sqrt{2}$.
	}
\end{ex}

%%=====Câu 11
\begin{ex}%[2D4V1-1]%[2015_K12 Huyên Nguyễn]
	Hàm số nào dưới đây có một nguyên hàm là $F(x)=\ln\left(x+\sqrt{1+x^2}\right)$?
	\choice
	{$f(x)=x\sqrt{1+x^2}$}
	{\True $f(x)=\dfrac{1}{\sqrt{1+x^2}}$}
	{$f(x)=\dfrac{x}{\sqrt{1+x^2}}$}
	{$f(x)=\sqrt{1+x^2}$}
	\loigiai{
		Ta có $F'(x)=\dfrac{\left(x+\sqrt{1+x^2}\right)'}{x+\sqrt{1+x^2}}=\dfrac{1+\dfrac{x}{\sqrt{x^2+1}}}{x+\sqrt{1+x^2}}=\dfrac{1}{\sqrt{x^2+1}}$.\\
		Vậy $F'(x)$ là một nguyên hàm của hàm số $f(x)=\dfrac{1}{\sqrt{1+x^2}}$.
	}
\end{ex}

%%=====Câu 12
\begin{ex}%[2D4V1-3]%[2015_K12 Huyên Nguyễn]
	Cho $F(x)$ và $G(x)$ lần lượt là nguyên hàm của hàm số $f(x)=\cos^2 x$ và $g(x)=\sin^2x$. Khi đó $F(x)-G(x)$ bằng
	\choice
	{$2\sin 2x+C$}
	{$2\cos 2x+C$}
	{$\cos 2x+C$}
	{\True $\dfrac{1}{2}\sin 2x+C$}
	\loigiai{
		Các hàm số $f(x)$ và $g(x)$ là các hàm số liên tục trên $\mathbb{R}$ nên ta có
		\begin{eqnarray*}
			F(x)+G(x)&=&\displaystyle\int f(x) \mathrm{\,d}x+\displaystyle\int g(x) \mathrm{\,d}x=\displaystyle\int \left[f(x)+g(x)\right] \mathrm{\,d}x\\
			&=&\displaystyle\int \left[\cos^2 x-\sin^2x\right] \mathrm{\,d}x=\displaystyle\int \cos 2x \mathrm{\,d}x\\
			&=&\dfrac{1}{2}\sin 2x+C.
		\end{eqnarray*}
	}
\end{ex}


\Closesolutionfile{ans}
\indapan{6}{ans/ans-2-B11-De2}

\cauds
\Opensolutionfile{ans}[ans/ans-2-B11-De2-DS]
\begin{ex}%[2D4H1-3]%[2015_K12 Huyên Nguyễn]
	Cho hàm số $f(x)=\left(\sin\dfrac{x}{2}+\cos \dfrac{x}{2}\right)^2$. 
	\choiceTF
	{$f(x)=1+\cos x$}
	{\True $f(x)$ liên tục trên $\mathbb{R}$}
	{$\displaystyle\int f(x)\mathrm{d}x=  x-\sin x +C$}
	{\True $F(x)$ là một nguyên hàm của $f(x)$ thỏa mãn $F(0)=1$. Khi đó $F(x)=x-\cos x+2$}
	\loigiai{
		\begin{itemchoice}
			\itemch Sai. $f(x)=\left(\sin\dfrac{x}{2}+\cos \dfrac{x}{2}\right)^2=\sin^2\dfrac{x}{2} +2\sin\dfrac{x}{2}\cos\dfrac{x}{2}+\cos^2\dfrac{x}{2}=1+\sin x$. 
			\itemch Đúng. $f(x)=1+\sin x$ liên tục trên $\mathbb{R}$.
			\itemch Sai. $\displaystyle\int f(x)\mathrm{d}x= \int \mathrm{d}x+\int \sin x \mathrm{d}x =x-\cos x +C$.
			\itemch Đúng. $F(x)=x-\cos x+C$ và $F(0)=1$ nên $0-\cos 0+C=1\Rightarrow C=2$.\\
			Vậy $F(x)=x-\cos x+2$
		\end{itemchoice}
	}
\end{ex}


%%Câu2
\begin{ex}%[2D4H1-4]%[2015_K12 Huyên Nguyễn]
	Cho $F(x)$ là một nguyên hàm của hàm số $f(x)=\mathrm{e}^x-\dfrac{4}{x}$ trong $\mathbb{R}\setminus \{0\}$ thỏa mãn $F(2)=\mathrm{e}^2$.
	\choiceTF
	{\True $F'(x)=\mathrm{e}^x-\dfrac{4}{x}$}
	{$f'(x)=\mathrm{e}^x+\ln|x|$}
	{$F(x)=\mathrm{e}^x-4\ln|x|-\mathrm{e}^2+4\ln2$}
	{Phương trình $F(x)=e^x$ có nghiệm duy nhất}
	\loigiai{
		\begin{itemchoice}
			\itemch Đúng. Vì $F(x)$ là một nguyên hàm của hàm số $f(x)$ nên $F'(x)=f(x)=\mathrm{e}^x-\dfrac{4}{x}$, $\forall x\ne 0$.
			\itemch Sai. Ta có $f'(x)=e^x+\dfrac{4}{x^2}$, $\forall x\ne 0$.
			\itemch Sai.  $F(x)=\displaystyle\int\left(\mathrm{e}^x-\dfrac{4}{x} \right)\mathrm{\,d}x=\mathrm{e}^x-4\ln|x|+C$.\\
			Lại có, $F(2)=\mathrm{e}^2\Leftrightarrow \mathrm{e}^2-4\ln2+C=\mathrm{e}^2\Leftrightarrow C=4\ln2$.\\
			Vậy $F(x)=\mathrm{e}^x-4\ln|x|+4\ln2$.
			\itemch Sai. Xét phương trình 
			\begin{eqnarray*}
				F(x)=\mathrm{e}^x &\Leftrightarrow&  \mathrm{e}^x-4\ln|x|+4\ln2 = \mathrm{e}^x\\
				&\Leftrightarrow& 	4\ln|x|=4\ln 2\\
				&\Leftrightarrow& \ln |x|=\ln 2\\
				&\Leftrightarrow&  |x|=2\\
				&\Leftrightarrow&\hoac{&x=-2\\&x=2.}
			\end{eqnarray*}
		\end{itemchoice}
	Vậy phương trình $F(x)=e^x$ có $2$ nghiệm.
	}
\end{ex}

%%Câu3
\begin{ex}%[2D4H1-2]%[2015_K12 Huyên Nguyễn]
	Cho hàm số $f(x)$ xác định trên $\mathbb{R}\backslash \left\{ \dfrac{1}{2}\right\}$, thỏa $f'\left( x \right) = \dfrac{2}{2x - 1}$, $f\left( 0 \right) = 1$ và $f\left( 1 \right) = 2$.
	\choiceTF
	{\True $f\left(x\right)=\ln\left| 2x-1\right|+C$ với $C\in \mathbb{R}$}
	{$f\left(x\right)=\ln\left(1-2x\right)+1\ \forall x\ne\dfrac{1}{2}$}
	{\True  Giá trị của biểu thức $f\left( { - 1} \right) + f\left( 3 \right)$ bằng $3+\ln 15$}
	{Với $x>\dfrac{1}{2}$ thì $\displaystyle\int e^{f(x)}\mathrm{d}x=2(x^2-x) +C$}
	\loigiai{
		\begin{itemchoice}
			\itemch Đúng. Ta có $f'\left(x\right)=\dfrac{2}{{2x-1}}$.\\
			Do đó, 
			$f\left(x\right)=\displaystyle\int\dfrac{2}{2x-1}\mathrm{d}x=\ln\left|2x-1\right|+C$ với $C\in \mathbb{R}$.
			\itemch Sai. $f(x)=\ln\left|2x-1\right|+C=\heva{& \ln\left({1-2x}\right)+{C_1}& \text{ với }x<\dfrac{1}{2}\\&\ln\left(2x-1\right)+C_2&\text{ với }x>\dfrac{1}{2}.}$\\
			Vì $f\left(0\right)=1\Rightarrow \ln\left({1-2.0}\right)+C_1=1\Rightarrow C_1=1$.\\
			$f\left(1\right)=2\Rightarrow \ln\left({2.1-1}\right)+C_2=2\Rightarrow C_2=2$.\\
			Do đó $f\left(x\right)=\heva{& \ln\left({1-2x}\right)+1& \text{ với }x<\dfrac{1}{2}\\&\ln\left(2x-1\right)+2&\text{ với }x>\dfrac{1}{2}.}$\\
			\itemch Đúng. Ta có $\heva{&f\left(-1\right)=\ln 3+1\\ &f\left(3\right)=\ln 5+2}\Rightarrow f\left(-1\right)+f\left(3\right)=3+\ln 5+\ln 3=3+\ln 15$.
			\itemch Sai. Với $x>\dfrac{1}{2}$ thì $f(x)=\ln (2x-1)+2\Rightarrow e^{f(x)}=e^2\cdot (2x-1)$.
			\begin{eqnarray*}
				\displaystyle\int e^{f(x)}\mathrm{d}x&=&\int e^2(2x-1)\mathrm{d}x\\
				&=&e^2\int (2x-1)\mathrm{d}x\\
				&=& e^2(x^2-x) +C.
			\end{eqnarray*}
		\end{itemchoice}
	}
\end{ex}

%%Câu 4
\begin{ex}%[2D4V1-6]%[2015_K12 Huyên Nguyễn]
	Một vận động viên điền kinh chạy với gia tốc $a(t)=-\dfrac{1}{24}t^3+\dfrac{5}{16}t^2$ (m/s$^2$), trong đó $t$ là khoảng thời gian tính từ lúc xuất phát.
	\choiceTF
	{ Vận tốc của vận động viên được tính theo công thức $v(t)=-\dfrac{1}{96}t^4+\dfrac{5}{48}t^3+3t$ m/s}
	{\True Vận tốc của vận động viên tại giây thứ $5$ là $6{,}51$ (m/s)}
	{Quãng đường chạy được của vận động viên sau thời gian $t$ giây được tính theo công thức $S(t)=-\dfrac{1}{480}t^5+\dfrac{5}{206}t^4$ (m)}
	{Sau $4$ giây thì vận động viên chạy được quãng đường là $6{,}2$ m}
	\loigiai{
		\begin{itemchoice}
			\itemch Sai.  $v(t)=\displaystyle\int a(t)\mathrm{\,d}t =\int \left( -\dfrac{1}{24}t^3+\dfrac{5}{16}t^2\right)\mathrm{\,d}t=-\dfrac{1}{96}t^4+\dfrac{5}{48}t^3+C$.\\
			Vì tại thời điểm $t=0$, vận tốc của vận động viên là $v(0)=0$ nên ta có\\ $v(t)=-\dfrac{1}{96}t^4+\dfrac{5}{48}t^3$ (m/s).
			\itemch Đúng. Tại giây thứ $5$ ta có 
			$v(5)=-\dfrac{1}{96}5^4+\dfrac{5}{48}5^3\approx 6{,}51$ m/s.
			\itemch Sai. Quãng đường chạy của vận động viên sau thời gian $t$ giây là 
			$$S(t)=\displaystyle\int v(t)\mathrm{\,d}t=\int \left( -\dfrac{1}{96}t^4+\dfrac{5}{48}t^3\right)\mathrm{\,d}t=-\dfrac{1}{480}t^5+\dfrac{5}{192}t^4+C.$$
			Tại thời điểm ban đầu $S(0)=0$ nên $S(t)=-\dfrac{1}{480}t^5+\dfrac{5}{192}t^4$ (m).
			\itemch Sai. Sau $4$ giây thì vận động viên chạy được quãng đường là\\ $S(4)=-\dfrac{1}{480}4^5+\dfrac{5}{192}4^4=\dfrac{68}{15}= 4{,}53$ m.
		\end{itemchoice}
	}
\end{ex}
\Closesolutionfile{ans}
\indapan{3}{ans/ans-2-B11-De2-DS}

\Opensolutionfile{ans}[ans/ans-2-B11-De2-KQ]
\caukq
\begin{ex}%[2D4H1-1]%[2015_K12 Huyên Nguyễn]
	Cho $F(x)=(ax+2b)\mathrm{e}^x$ ($a,b\in\mathbb{R}$) là một nguyên hàm của hàm số $f(x)=(x+5)\mathrm{e}^x$. Tính giá trị của $a^2+2b^2$.
	\shortans{$9$}
	\loigiai{
		Vì $F(x)$ là một nguyên hàm của hàm số $f(x)$ nên 
		$$F'(x)=f(x)\Leftrightarrow \left[ax+(a+2b)\right]\mathrm{e}^x=(x+5)\mathrm{e}^x\Leftrightarrow \heva{& a=1 \\ & a+2b=5}\Leftrightarrow \heva{& a=1 \\ & b=2.}$$
		Vậy $a^2+2b^2=9$.
	}
\end{ex}
\begin{ex}%[2D4H1-2]%[2015_K12 Huyên Nguyễn]
	Cho hàm số $f(x)$ xác định trên $\mathbb{R}\setminus\{1\}$ thỏa mãn $f'(x)=\dfrac{1}{x-1}$, $f(0)=2$, $f(4)=3$. Tính giá trị của biểu thức $f(-2)+f(2)$.
		\shortans{$5$}
	\loigiai{
		Ta có $f(x)=\displaystyle\int f'(x) \mathrm{\,d}x=\displaystyle\int \dfrac{1}{x-1}\mathrm{\,d}x=\ln|x-1|+C=\heva{& \ln(x-1)+C_1&\text{ với }x> 1 \\ & \ln (1-x)+C_2&\text{ với }x<1.}$\\
		Vì $f(0)=2\Rightarrow \ln (1-0)+C_2=2\Rightarrow C_2=2$;\\ $f(4)=3\Rightarrow \ln (4-1)+C_1=3\Rightarrow C_1=3-\ln 3$.\\
		Suy ra $f(x)=\heva{& \ln(x-1)+3-\ln 3&\text{ với }x> 1 \\ & \ln (1-x)+2&\text{ với }x<1.}$\\
		Suy ra $\heva{&f\left(-2\right)=\ln 3+2\\ &f\left(2\right)=3-\ln 3}\Rightarrow f\left(-2\right)+f\left(2\right)=5$.
	}
\end{ex}

%%==========Câu 3
\begin{ex}%[2D4V1-6]
	Một vật chuyển động có gia tốc là $a\left(t\right)=\dfrac{3}{t+1}$ (m/s$^2$). Biết rằng vận tốc ban đầu của vật là $6$ m/s. Vận tốc của vật đó sau $5$ giây là bao nhiêu m/s$^2$ (làm tròn kết quả đến hàng phần mười)?
	\shortans{$13{,}2$}
	\loigiai{
		Vận tốc của vật là $v(t)=\displaystyle \int a\left(t\right)  \mathrm{\,d}t =\int \dfrac{3}{t+1}\mathrm{\,d}t=3\ln |t+1|+C$.\\
		Vì vận tốc tốc ban đầu của vật là $6$ m/s nên $v(0)=6 \Leftrightarrow 3\ln 1+C=6  \Leftrightarrow C=6$.\\
		Suy ra $v(t)=3\ln |t+1| +6$.\\
		Vận tốc của vật đó sau $10$ giây là $v(10)=3\ln|10+1|+6\approx 13{,}2$ (m/s).
	}
\end{ex}

%%==========Câu 4
\begin{ex}%[2D4V1-1]%[2015_K12 Huyên Nguyễn]
	Biết $\displaystyle\int   \dfrac{2x+3}{\mathrm{e}^x}\text{d}x= (ax+b)\mathrm{e}^{-x}+C$ (với $a$, $b\in \mathbb{R}$). Tính giá trị của $a-2b$.
	\shortans{$8$}
	\loigiai
	{		
		Đặt $F(x)=(ax+b)\mathrm{e}^{-x}+C\Rightarrow F'(x)=\dfrac{2x+3}{\mathrm{e}^x}=(2x+3)\mathrm{e}^{-x}$.
		\begin{eqnarray*}
			F'(x)=(2x+3)\mathrm{e}^{-x} &\Leftrightarrow& -(ax+b)\mathrm{e}^{-x}+a  \mathrm{e}^{-x}=(2x+3)\mathrm{e}^{-x} \\
			&\Leftrightarrow&(-ax+a-b)\mathrm{e}^{-x}=(2x+3)\mathrm{e}^{-x}\\
			&\Leftrightarrow &\heva{&-a=2\\& a-b=3.}
		\end{eqnarray*}
		Suy ra $\heva{&a=-2\\&b=-5}\Rightarrow a-2b =8$.
	} 
\end{ex}
%%==========Câu 5
\begin{ex}%[2D4V1-3]%[2015_K12 Huyên Nguyễn]
	Cho $F(x)$ là một nguyên hàm của hàm số $f(x)=\cos 2x$ và thỏa mãn $F(\pi)=1$. Phương trình $F(x)=1$ có tất cả bao nhiêu nghiệm trong đoạn $[0 ; 3 \pi]$?
	\shortans{$7$}	
\loigiai{
	Ta có $F(x)=\displaystyle\int\cos 2x\mathrm{\,d}x=\dfrac{1}{2}\sin 2x+C$.\\
	$F(\pi)=1\Leftrightarrow C=1$, suy ra $F(x)=\dfrac{1}{2}\sin 2x +1$.\\
	Phương trình
	\begin{eqnarray*}
		F(x)=1&\Leftrightarrow&\dfrac{1}{2}\sin 2x +1=1\\
		&\Leftrightarrow&\sin 2x=0\\
		&\Leftrightarrow& 2x=k\pi,\ k\in \mathbb{Z}\\
		&\Leftrightarrow& x=\dfrac{k\pi}{2},\ k\in \mathbb{Z}
	\end{eqnarray*}
	Vậy trong đoạn $[0;3\pi]$ phương trình có tất cả $7$ nghiệm $x\in \left\{0;\dfrac{\pi}{2};\pi; \dfrac{3\pi}{2};2\pi; \dfrac{5\pi}{2};3\pi\right\}$.
	}
\end{ex}


%%==========Câu 6
\begin{ex}%[2D4V1-2]%[2015_K12 Huyên Nguyễn]
	Cho $F(x)$ là một nguyên hàm của hàm số $f(x)=\dfrac{1}{x}$ và thỏa mãn $F(1)=-1$. Có tất cả bao nhiêu giá trị của $x$ để $2F(x)=\dfrac{1}{F(x)+1}-1 $?
	\shortans{$4$}
	\loigiai{
		Ta có $F(x)=\displaystyle\int\dfrac{1}{x}\mathrm{\,d}x=\ln|x|+C$.\\
		Mà $F(1)=-1\Leftrightarrow \ln|1|+C=-1\Leftrightarrow C=-1$, suy ra $F(x)=\ln |x|-1$.\\
		Do đó 
		\begin{eqnarray*}
			2F(x)=\dfrac{1}{F(x)+1}-1&\Leftrightarrow& 2\ln |x|-1=\dfrac{1}{\ln|x|}\\
			&\Leftrightarrow &2\ln^2|x|-\ln|x|-1=0\\
			&\Leftrightarrow & \hoac{&\ln|x|=1\\&\ln|x|=-\dfrac{1}{2}}\\
			&\Leftrightarrow & x\in \left\{\pm \mathrm{e};\pm \dfrac{1}{\sqrt{\mathrm{e}}}\right\}.	
		\end{eqnarray*}
		Vậy có $4$ giá trị của $x$ để $2F(x)=\dfrac{1}{F(x)+1}-1 $.
	}
\end{ex}


\Closesolutionfile{ans}
\indapan{6}{ans/ans-2-B11-De2-KQ}